\newglossaryentry{s3}{
    name={S3},
    description={протокол объектного хранилища, изначально предложенный Amazon Simple Storage Service и ставший де-факто стандартом}
}

\newglossaryentry{staas}{
    name={STaaS},
    description={(\textit{Storage as a Service})~--- модель предоставления хранилищ как сервиса; внутренняя платформа <<Авито>>, объединяющая S3-кластеры и предоставляющая командам доступ к ним по единому API}
}

\newglossaryentry{bucket}{
    name={Bucket},
    description={контейнер для хранения объектов в S3, имеющий уникальное имя в пределах хранилища}
}

\newglossaryentry{rclone}{
    name={RClone},
    description={свободная утилита для синхронизации данных между файловыми системами и облачными хранилищами}
}

\newglossaryentry{consul}{
    name={Consul},
    description={сервис обнаружения, конфигурации и распределённой координации от HashiCorp; в \gls{staas} используется как DNS-сервис для маршрутизации клиентов к S3-кластерам и как механизм распределённых блокировок}
}

\newglossaryentry{vault}{
    name={Vault},
    description={хранилище секретов от HashiCorp; в \gls{staas} хранит учётные данные доступа к S3-бакетам}
}

\newglossaryentry{slo}{
    name={SLO},
    description={(\textit{Service Level Objective})~--- целевой показатель качества обслуживания, который платформа гарантирует клиентам, например 95-й перцентиль времени ответа API}
}
